% ED26-SRC-063-0103: chapter pp.101--246; this packet treats §§1--2.2, pp.101--165.
\chapter{Aristotle's Underlying Logic: Boger's Historical Reconstruction}
\label{sec:p13-src063-0103}
\chaptermark{Aristotle's Underlying Logic}

George Boger's 2004 chapter distinguishes Aristotle's own investigations
from later Aristotelian logic. Its object is an \emph{underlying logic}: a
language, its semantics, and an apparatus for deduction. It connects
\emph{Categories}, \emph{On Interpretation}, \emph{Prior Analytics},
\emph{Posterior Analytics}, \emph{Metaphysics}, \emph{Topics}, and
\emph{Sophistical Refutations}; a few syllogistic moods are not the whole
ancient theory \parencite[pp.~101--118]{Boger2004UnderlyingLogic}.
The claim that Aristotle intentionally modeled such a logic is Boger's
historical interpretation. The organization into syntax, semantics and
proof theory is his explanatory template, not a verbatim ancient textbook.

\section{Scientific purpose and competing interpretations}
\label{sec:p13-boger-purpose}

Boger locates Aristotle within a fifth- and fourth-century movement toward
naturalistic explanation, observation, argument, and scientific organization.
He connects Protagoras's political teaching with Hippocratic medicine, rather
than treating ancient rationality as exclusively mathematical. Aristotle's
ideal includes scientific and deliberative excellence. Boger also notices
the social hierarchy in the passage about master workers and manual
laborers; the historical narrative does not validate every associated social
judgment \parencite[pp.~101--104]{Boger2004UnderlyingLogic}.

The medical example in \emph{Metaphysics} 981a5--981b6 separates experience
of particular successful treatments from an art grasping a general relation
and its explanation. Knowing the art without knowing the individual patient
can still produce failure. Repeated practical success, conversely, is not
by itself knowing why. Logic has an epistemic purpose without all knowledge
being identified with deduction.

Boger distinguishes three interpretive families. Traditional presentations
often conflate a fully interpreted premise--conclusion argument with an
uninterpreted pattern, and test it through quantity, quality, and
distribution. For Aristotle, in Boger's reading, an invalid syllogism is a
contradiction in terms: a syllogism already has necessary consequence.
An invalid candidate argument is not therefore a syllogism.

The axiomaticist interpretation of \L{}ukasiewicz, Miller, Bochenski, and
Patzig identifies syllogisms with logically true conditionals, some taken as
axioms, and emphasizes formal ontology. The deductionist interpretation
of Corcoran, Smiley, and Smith emphasizes cogent deductions, possibly with
more than two premises, and natural deduction or formal epistemology.
Boger does not simply adopt either. Aristotle reasoned in an object
language and studied that reasoning in a metalanguage. Insufficient
separation of these levels, Boger argues, obscures the ancient project
\parencite[pp.~104--107]{Boger2004UnderlyingLogic}.

Two disagreements matter for lookup. Boger rejects the common
identification of \emph{deduction}, \emph{reduction}, and \emph{analysis}.
He also denies that the contrasted-instance procedure of
\emph{Prior Analytics} A4--6 should simply be called a modern
counterargument or counterinterpretation method. Johnson's different
countermodel terminology must not replace this source-specific position.

\section{Deduction is not demonstration}
\label{sec:p13-boger-demonstration}

Boger reads \emph{Prior Analytics} 25b26--31 as investigating how syllogisms
arise, and \emph{Posterior Analytics} 71b17--25 as specifying demonstrative
understanding. Deduction establishes that a conclusion follows necessarily.
It need not establish that the premises are true, and can use false premises.
A demonstration is a special deduction whose premises are true, primitive,
immediate, better known, prior, and explanatory of its conclusion.
Calling a valid derivation scientific does not supply these further
conditions \parencite[pp.~107--110]{Boger2004UnderlyingLogic}.

An ontic fact is distinct from an epistemic achievement. Truth and
necessary consequence are not made to hold by a participant's recognition;
a participant may nevertheless fail to know them. A science's first
principles must be established independently of the deduction that uses
them. Derived theorems may become premises later in the extended discourse.
True premises and validity secure a true conclusion, but do not by
themselves secure every explanatory condition of demonstration.

Boger compares common notions in mathematics with elementary categorical
deduction patterns. Their generality is relative to the relevant kinds
of being, not complete absence of interpretation. Reflection on axioms
and noncontradiction belongs to philosophy's study of being as being.
This supports the proposed rule-nature of syllogistic patterns without
showing that Aristotle performed a complete metalogical study of
geometric proof \parencite[pp.~109--114]{Boger2004UnderlyingLogic}.

\begin{hbdefinition}{def:p13-boger-discourse-objects}{Distinct logical objects}
A sentence is true or false. An argument comprises premises and a conclusion
and is valid or invalid. An argumentation additionally has a chain of
reasoning; a cogent one is a deduction. An argument pattern is
metalinguistic. Boger calls it \emph{panvalid}, \emph{paninvalid}, or
\emph{neutrovalid} according as its instances are all valid, all invalid,
or some of each. These classify instances, not truth of schematic letters.
\end{hbdefinition}

A concludent categorical premise-pair pattern supports a necessarily
following categorical result; an inconcludent one does not. Traditional
names such as Barbara name patterns and, where appropriate, rules, not
individual arguments. Boger's narrower rendering of \emph{sullogismos}
in A4--7 uses a valid two-premise, three-term argument in one of three
figures; other contexts require the broader rendering deduction.
Note~13 records his change from identifying the word with the pattern
itself \parencite[pp.~115--117]{Boger2004UnderlyingLogic}.

\section{Natural language, declarative sentences, and meaning}
\label{sec:p13-boger-natural-grammar}

Boger begins sentence formation with combining a noun and a verb or predicate
into meaningful speech. Meaningful words placed side by side do not
automatically form a unified sentence; a noun alone affirms nothing.
A declarative sentence is the special kind that can be true or false.
Prayers and other nondeclarative uses are not truth-bearing premises
merely because they have grammatical structure.

His extracted SFR1--SFR3 concern meaningful combination, simple predication
of one thing of one thing, and contradiction by negating the whole
affirmation. They are reconstructed principles, not a complete recursive
grammar for Greek. Boger acknowledges mixed syntactic/semantic vocabulary
and lack of modern rigor \parencite[pp.~118--126]{Boger2004UnderlyingLogic}.

A sentence, its meaning, and its denoted state of affairs are different
objects. Different expressions may convey the same meaning; one may
convey several. Assigning the name garment to both a horse and a man
makes the apparent sentence about a white garment fail to be one simple
affirmation. This supports the distinction without establishing that
Aristotle possessed modern model theory.

A sentence denial differs from a negative term.
\enquote{Socrates is not good} denies an affirmation;
\enquote{Socrates is not-good} can affirm a property expressed by an
indefinite predicate. A contrary predicate such as bad is different again.
A negative-looking expression is insufficient to identify a contradictory
\parencite[pp.~124--128, 131--136]{Boger2004UnderlyingLogic}.

\section{The artificial categorical language}
\label{sec:p13-boger-categorical-grammar}

Boger writes \(AaB,AeB,AiB,AoB\). Here \(A\) is the predicate placeholder,
\(B\) the subject placeholder, and the lower-case middle letter represents
belonging. This differs from Johnson's subject-first \(Aab,Eab,Iab,Oab\).

\begin{tabularx}{\textwidth}{@{}l X X@{}}
\toprule
Pattern & Predication & Quality and quantity\\
\midrule
\(AaB\) & \(A\) belongs to every \(B\). & Universal attributive\\
\(AeB\) & \(A\) belongs to no \(B\). & Universal privative\\
\(AiB\) & \(A\) belongs to some \(B\). & Partial attributive\\
\(AoB\) & \(A\) does not belong to some \(B\), or not to every \(B\). & Partial privative\\
\bottomrule
\end{tabularx}

CSFR1 fixes predicate, logical constant, then subject. CSFR3 takes the
two nonlogical expressions as homogeneous substantives, making conversion
intelligible. CSFR2 restricts the scientific presentation to different
terms, but note~29 acknowledges Aristotle's identity and circular-reasoning
cases; Aristotle could understand \enquote{Every horse is a horse}.
CSFR4 closes this formation account. These restrictions are not a grammar
for all Greek sentences \parencite[pp.~128--133]{Boger2004UnderlyingLogic}.

Uppercase letters are schematic placeholders, not bound/free variables
ranging over individuals. Aristotle supplied no abbreviations for the
four logical constants. Actual terms belong to interpreted topical
discourse; their places belong to the metalanguage. An unfilled pattern
has no truth value. This is not a wholly uninterpreted modern calculus,
since logical constants retain intended meanings
\parencite[pp.~141--144]{Boger2004UnderlyingLogic}.

Contradictories keep both term places and exchange \(a/o\) or \(e/i\).
Universal contraries exchange \(a/e\). Contradiction excludes both joint
truth and joint falsity; contrariety excludes joint truth but permits
joint falsity in the intended existent-subject cases.
Partial attribution and partial privation are not contradictory simply
because one contains a negative operator
\parencite[pp.~133--136, 155--157]{Boger2004UnderlyingLogic}.

The formalization also serves reasoning practice. Boger contrasts
structural training with memorizing stock speeches. Analysis transforms
between figures while preserving the problem proved. A topical science has
a genus, vocabulary, principles, and attributes. Cogent reasoning may
nevertheless be irrelevant to it: Zeno's general argument is not a medical
objection to walking after dinner. Optics/geometry illustrate qualified
crossing through subordination, not arbitrary transfer of demonstrations
\parencite[pp.~136--142]{Boger2004UnderlyingLogic}.

\section{Truth, existence, and an explicitly limited bridge}
\label{sec:p13-boger-semantics}

Attribution in things differs from predication by a participant.
A sentence is true when it accords with the relevant state of affairs.
Uttering it does not make that fact obtain. Someone sitting and then rising
illustrates how an unchanged expression changes truth with its subject
matter. Truth/falsity of sentences differ from health/sickness or
sight/blindness as qualities and privations of substances.
Boger reads Aristotle's use of truth vocabulary for both sentences and
what obtains as context-sensitive, not automatically a category mistake
\parencite[pp.~144--157]{Boger2004UnderlyingLogic}.
His correspondence comparison does not attribute Tarski's full formal
definition to Aristotle.

Boger distinguishes two readings of some. A determinate reading says
some members have the property and some lack it. The indeterminate reading
says at least one has it, perhaps all. General expressions such as
\enquote{Men are white} can leave the scope unsettled. Rejection arguments
sometimes use a universal denial to establish an indeterminate partial
denial; this implication is not justified under an unannounced
\enquote{some but not all} replacement
\parencite[pp.~147--148]{Boger2004UnderlyingLogic}.

The existential discussion is not the modern convention that every
universal is vacuously true of an empty class. In Boger's tables~7--9,
affirmations about a nonexistent subject are false and genuine denials
true. A positive contrary predicate is another affirmation and also false.
With an existent subject, contrary predicates need not exhaust the
possibilities: sight/blindness requires the relevant natural condition and
time. A contradictory is instead the denial of the very same affirmation
\parencite[pp.~157--164]{Boger2004UnderlyingLogic}.

Thus \enquote{Socrates ails} and \enquote{Socrates fares} may both be false
for a nonexistent subject, whereas \enquote{Socrates does not ail} is the
true denial in this interpretation. Universally, both
\enquote{Every man is good} and \enquote{Every man is bad} are affirmations;
\enquote{No man is good} is a denial. Partially, \enquote{Some man is good}
and \enquote{Some man is bad} are affirmations, whereas
\enquote{Some man is not good} contains the negative operator.
Boger identifies table~9's subcontrary organization as his interpolation,
not an ancient technical term reproduced verbatim.

\begin{hbdefinition}{def:p13-boger-guarded-bridge}{Modern comparison semantics, not historical ontology}
Let \(S,P\) be subject and predicate extensions in a nonempty domain.
For a local comparison of the nonexistence policy define
\[
 \alpha(S,P)\mathrel{:\equiv} S\ne\varnothing\land S\subseteq P,\qquad
 \iota(S,P)\mathrel{:\equiv} S\cap P\ne\varnothing,
\]
\[
 \varepsilon(S,P)\mathrel{:\equiv}\neg\iota(S,P),\qquad
 \omega(S,P)\mathrel{:\equiv}\neg\alpha(S,P).
\]
They correspond locally to \(AaB,AiB,AeB,AoB\), respectively.
This extensional comparison does not reproduce every intensional feature.
For determinate some-but-not-all one instead requires both
\(S\cap P\ne\varnothing\) and \(S\setminus P\ne\varnothing\).
\end{hbdefinition}

\begin{hbtheorem}{thm:p13-boger-existence-bridge}{proved=annex:p13.boger.existence}{What nonemptiness changes}
For \(S\ne\varnothing\), these conditions agree with the ordinary
extensional universal, particular, universal-negative, and
particular-negative conditions. The pairs \(\alpha/\omega\) and
\(\varepsilon/\iota\) are contradictory for every \(S\).
On nonempty \(S\), \(\alpha,\varepsilon\) cannot both hold but may both
fail; \(\iota,\omega\) cannot both fail but may both hold.
For \(S=\varnothing\), the values of
\(\alpha,\iota,\varepsilon,\omega\) are false, false, true, true.
Ordinary vacuous universality instead makes \(S\subseteq P\) true and
\(S\setminus P\ne\varnothing\) false.
\end{hbtheorem}

A two-member subject with only one member in \(P\) realizes both false
universals and both true indeterminate particulars. An empty subject shows
why an unqualified lookup replacement by set inclusion changes the
interpretation; it does not establish this as the only historical reading.

Boger prefers an intensional or possessional account: an individual has
or lacks a property, and essential versus accidental attribution matters.
Modern class inclusion is useful comparison, not evidence that Aristotle
identified substance with a set. Meaning, predication, and existential
import must accompany the formula \parencite[pp.~164--165]{Boger2004UnderlyingLogic}.
